\PassOptionsToPackage{unicode}{hyperref}
\PassOptionsToPackage{hyphens}{url}
\PassOptionsToPackage{dvipsnames,svgnames,x11names}{xcolor}
%
\documentclass[
  11pt,
  letterpaper,
  DIV=11,
  numbers=noendperiod]{scrartcl}
\usepackage{xcolor}
\usepackage[margin=1in]{geometry}
\usepackage{amsmath,amssymb}
\setcounter{secnumdepth}{-\maxdimen} % remove section numbering
\usepackage{iftex}
\ifPDFTeX
  \usepackage[T1]{fontenc}
  \usepackage[utf8]{inputenc}
  \usepackage{textcomp}
\else
  \usepackage{unicode-math}
  \defaultfontfeatures{Scale=MatchLowercase}
  \defaultfontfeatures[\rmfamily]{Ligatures=TeX,Scale=1}
\fi
\usepackage{lmodern}
\IfFileExists{upquote.sty}{\usepackage{upquote}}{}
\IfFileExists{microtype.sty}{%
  \usepackage[]{microtype}
  \UseMicrotypeSet[protrusion]{basicmath}
}{}
\makeatletter
\@ifundefined{KOMAClassName}{%
  \IfFileExists{parskip.sty}{%
    \usepackage{parskip}
  }{%
    \setlength{\parindent}{0pt}
    \setlength{\parskip}{6pt plus 2pt minus 1pt}}
}{%
  \KOMAoptions{parskip=half}}
\makeatother

% Code listings (verbatim with light shading)
\usepackage{color}
\usepackage{fancyvrb}
\newcommand{\VerbBar}{|}
\newcommand{\VERB}{\Verb[commandchars=\\\{\}]}
\DefineVerbatimEnvironment{Highlighting}{Verbatim}{commandchars=\\\{\}}
\usepackage{framed}
\definecolor{shadecolor}{RGB}{241,243,245}
\newenvironment{Shaded}{\begin{snugshade}}{\end{snugshade}}
\newcommand{\AlertTok}[1]{\textcolor[rgb]{0.68,0.00,0.00}{#1}}
\newcommand{\AnnotationTok}[1]{\textcolor[rgb]{0.37,0.37,0.37}{#1}}
\newcommand{\AttributeTok}[1]{\textcolor[rgb]{0.40,0.45,0.13}{#1}}
\newcommand{\BaseNTok}[1]{\textcolor[rgb]{0.68,0.00,0.00}{#1}}
\newcommand{\BuiltInTok}[1]{\textcolor[rgb]{0.00,0.23,0.31}{#1}}
\newcommand{\CharTok}[1]{\textcolor[rgb]{0.13,0.47,0.30}{#1}}
\newcommand{\CommentTok}[1]{\textcolor[rgb]{0.37,0.37,0.37}{#1}}
\newcommand{\CommentVarTok}[1]{\textcolor[rgb]{0.37,0.37,0.37}{\textit{#1}}}
\newcommand{\ConstantTok}[1]{\textcolor[rgb]{0.56,0.35,0.01}{#1}}
\newcommand{\ControlFlowTok}[1]{\textcolor[rgb]{0.00,0.23,0.31}{\textbf{#1}}}
\newcommand{\DataTypeTok}[1]{\textcolor[rgb]{0.68,0.00,0.00}{#1}}
\newcommand{\DecValTok}[1]{\textcolor[rgb]{0.68,0.00,0.00}{#1}}
\newcommand{\DocumentationTok}[1]{\textcolor[rgb]{0.37,0.37,0.37}{\textit{#1}}}
\newcommand{\ErrorTok}[1]{\textcolor[rgb]{0.68,0.00,0.00}{#1}}
\newcommand{\ExtensionTok}[1]{\textcolor[rgb]{0.00,0.23,0.31}{#1}}
\newcommand{\FloatTok}[1]{\textcolor[rgb]{0.68,0.00,0.00}{#1}}
\newcommand{\FunctionTok}[1]{\textcolor[rgb]{0.28,0.35,0.67}{#1}}
\newcommand{\ImportTok}[1]{\textcolor[rgb]{0.00,0.46,0.62}{#1}}
\newcommand{\InformationTok}[1]{\textcolor[rgb]{0.37,0.37,0.37}{#1}}
\newcommand{\KeywordTok}[1]{\textcolor[rgb]{0.00,0.23,0.31}{\textbf{#1}}}
\newcommand{\NormalTok}[1]{\textcolor[rgb]{0.00,0.23,0.31}{#1}}
\newcommand{\OperatorTok}[1]{\textcolor[rgb]{0.37,0.37,0.37}{#1}}
\newcommand{\OtherTok}[1]{\textcolor[rgb]{0.00,0.23,0.31}{#1}}
\newcommand{\PreprocessorTok}[1]{\textcolor[rgb]{0.68,0.00,0.00}{#1}}
\newcommand{\RegionMarkerTok}[1]{\textcolor[rgb]{0.00,0.23,0.31}{#1}}
\newcommand{\SpecialCharTok}[1]{\textcolor[rgb]{0.37,0.37,0.37}{#1}}
\newcommand{\SpecialStringTok}[1]{\textcolor[rgb]{0.13,0.47,0.30}{#1}}
\newcommand{\StringTok}[1]{\textcolor[rgb]{0.13,0.47,0.30}{#1}}
\newcommand{\VariableTok}[1]{\textcolor[rgb]{0.07,0.07,0.07}{#1}}
\newcommand{\VerbatimStringTok}[1]{\textcolor[rgb]{0.13,0.47,0.30}{#1}}
\newcommand{\WarningTok}[1]{\textcolor[rgb]{0.37,0.37,0.37}{\textit{#1}}}

\usepackage{enumitem}
\setlist[enumerate]{itemsep=10pt}

% Callout box (note style)
\usepackage[skins,breakable]{tcolorbox}
\usepackage{fontawesome5}
\definecolor{quarto-callout-note-color}{HTML}{0758E5}
\definecolor{quarto-callout-note-color-frame}{HTML}{4582ec}

\setlength{\emergencystretch}{3em}
\providecommand{\tightlist}{%
  \setlength{\itemsep}{0pt}\setlength{\parskip}{0pt}}

\usepackage{bookmark}
\IfFileExists{xurl.sty}{\usepackage{xurl}}{}
\urlstyle{same}
\hypersetup{
  pdftitle={Quiz 1 --- Practice},
  colorlinks=true,
  linkcolor={blue},
  filecolor={Maroon},
  citecolor={Blue},
  urlcolor={Blue},
  pdfcreator={LaTeX}}

\title{Quiz 1 --- Practice}
\usepackage{etoolbox}
\makeatletter
\providecommand{\subtitle}[1]{%
  \apptocmd{\@title}{\par {\large #1 \par}}{}{}
}
\makeatother
\subtitle{STAT 184 --- Summer 2026}
\author{}
\date{}

\begin{document}
\maketitle

\noindent \textbf{Name:} \rule{2in}{0.4pt} \hfill \textbf{Date:} \rule{2in}{0.4pt}

\vspace{0.5em}

\begin{tcolorbox}[enhanced jigsaw, arc=.35mm, bottomrule=.15mm, bottomtitle=1mm, breakable, colback=white, colbacktitle=quarto-callout-note-color!10!white, colframe=quarto-callout-note-color-frame, coltitle=black, left=2mm, leftrule=.75mm, opacityback=0, opacitybacktitle=0.6, rightrule=.15mm, title=\textcolor{quarto-callout-note-color}{\faInfo}\hspace{0.5em}{About this practice quiz}, titlerule=0mm, toprule=.15mm, toptitle=1mm]

\begin{itemize}
\tightlist
\item This is a \textbf{practice} quiz to help you prepare for Quiz 1.
\item It mirrors the \textbf{format, length, and topic distribution} of the real quiz.
\item The actual quiz will have different specific questions, but they will test the same concepts.
\item \textbf{Recommended setup:} Give yourself 20 minutes, closed-book, no electronics. Then check your answers against the key.
\item 9 questions, 30 points total.
\end{itemize}

\end{tcolorbox}

\vspace{0.5em}

\subsection{1. (2 pts) Multiple choice --- assignment}

Which of the following is the \textbf{preferred} way (in this course) to assign the value \texttt{100} to an object named \texttt{score}?

\vspace{0.25em}

\begin{enumerate}
\def\labelenumi{(\Alph{enumi})}
\tightlist
\item \texttt{score\ =\ 100}
\item \texttt{score\ ==\ 100}
\item \texttt{score\ \textless{}-\ 100}
\item \texttt{100\ \textless{}-\ score}
\item \texttt{assign(score,\ 100)}
\end{enumerate}

\vspace{0.25em}

$\textbf{Answer:}$ \rule{1in}{0.4pt}

\newpage

\subsection{2. (2 pts) Multiple choice --- packages}

You just opened RStudio for the day and want to use functions from the \texttt{ggplot2} package. You installed \texttt{ggplot2} a few weeks ago. Which command do you need to run to make the functions available in your current session?

\vspace{0.25em}

\begin{enumerate}
\def\labelenumi{(\Alph{enumi})}
\tightlist
\item Reinstall the package with \texttt{install.packages("ggplot2")}.
\item Load it with \texttt{library(ggplot2)}.
\item Nothing --- once installed, it's loaded automatically every session.
\item Restart R and try again.
\item Run \texttt{update.packages("ggplot2")} first.
\end{enumerate}

$\textbf{Answer:}$ \rule{1in}{0.4pt}

\subsection{3. (4 pts) Code reading --- predict the output}

For each of the following, write what R prints to the console.

\vspace{0.25em}

\begin{enumerate}
\def\labelenumi{(\alph{enumi})}
\tightlist
\item \texttt{sqrt(144)}
\end{enumerate}

\begin{quote}
\rule{4in}{0.4pt}
\end{quote}

\vspace{0.5em}

\begin{enumerate}
\def\labelenumi{(\alph{enumi})}
\setcounter{enumi}{1}
\tightlist
\item \texttt{c(10,\ 20,\ 30)\ +\ 5}
\end{enumerate}

\begin{quote}
\rule{4in}{0.4pt}
\end{quote}

\vspace{0.5em}

\begin{enumerate}
\def\labelenumi{(\alph{enumi})}
\setcounter{enumi}{2}
\tightlist
\item \texttt{15\ \%\%\ 4}
\end{enumerate}

\begin{quote}
\rule{4in}{0.4pt}
\end{quote}

\vspace{0.5em}

\begin{enumerate}
\def\labelenumi{(\alph{enumi})}
\setcounter{enumi}{3}
\tightlist
\item \texttt{mean(c(1,\ 3,\ 5,\ 7,\ 9))}
\end{enumerate}

\begin{quote}
\rule{4in}{0.4pt}
\end{quote}

\vspace{0.5em}

\newpage

\subsection{4. (4 pts) Code reading --- identify the parts}

Consider the following code:

\begin{Shaded}
\begin{Highlighting}[]
\NormalTok{Movies }\SpecialCharTok{|\textgreater{}}
  \FunctionTok{filter}\NormalTok{(year }\SpecialCharTok{==} \DecValTok{2010}\NormalTok{) }\SpecialCharTok{|\textgreater{}}
  \FunctionTok{summarise}\NormalTok{(}\AttributeTok{avg\_rating =} \FunctionTok{mean}\NormalTok{(rating))}
\end{Highlighting}
\end{Shaded}

\vspace{0.25em}

\begin{enumerate}
\def\labelenumi{(\alph{enumi})}
\item Name the \textbf{data table}:

  \rule{2in}{0.4pt}
\item Name \textbf{two functions} used:

  \rule{2.5in}{0.4pt}
\item Name \textbf{one variable} (column):

  \rule{2in}{0.4pt}
\item Name \textbf{one constant}:

  \rule{2in}{0.4pt}
\end{enumerate}

\vspace{0.5em}

\subsection{5. (3 pts) Multiple choice --- data structures}

Which of the following best describes a \textbf{vector} in R?

\vspace{0.25em}

\begin{enumerate}
\def\labelenumi{(\Alph{enumi})}
\tightlist
\item A one-dimensional collection of values, all of the same type.
\item A two-dimensional collection of values, all of the same type.
\item A two-dimensional collection where each column can hold a different type.
\item An ordered list of unrelated R objects, possibly of different types.
\item A single number stored in an object.
\end{enumerate}

\vspace{0.25em}

$\textbf{Answer:}$ \rule{1in}{0.4pt}

\newpage

\subsection{6. (3 pts) Code reading --- function arguments}

Consider this call:

\begin{Shaded}
\begin{Highlighting}[]
\FunctionTok{round}\NormalTok{(}\FloatTok{3.14159}\NormalTok{, }\AttributeTok{digits =} \DecValTok{3}\NormalTok{)}
\end{Highlighting}
\end{Shaded}

\vspace{0.25em}

\begin{enumerate}
\def\labelenumi{(\alph{enumi})}
\item What is the \textbf{name of the function} being called?

  \rule{2in}{0.4pt}
\item How many \textbf{arguments} are passed to it?

  \rule{1in}{0.4pt}
\item Is the first argument (\texttt{3.14159}) \textbf{named} or \textbf{positional}?

  \rule{1.5in}{0.4pt}
\end{enumerate}

\vspace{0.5em}

\subsection{7. (3 pts) Short answer --- comments on code}

In one or two sentences, explain the difference between running \texttt{mean(c(2,\ 4,\ 6))} and running \texttt{m\ \textless{}-\ mean(c(2,\ 4,\ 6))}. Which one stores the result for later use?

\vspace{0.5em}

\rule{6.5in}{0.4pt}

\vspace{1em}

\rule{6.5in}{0.4pt}

\vspace{1em}

\rule{6.5in}{0.4pt}

\vspace{0.5em}

\newpage

\subsection{8. (4 pts) Code reading --- histograms}

Consider this code:

\begin{Shaded}
\begin{Highlighting}[]
\FunctionTok{ggplot}\NormalTok{(penguins, }\FunctionTok{aes}\NormalTok{(}\AttributeTok{x =}\NormalTok{ bill\_length\_mm)) }\SpecialCharTok{+}
  \FunctionTok{geom\_histogram}\NormalTok{(}\AttributeTok{binwidth =} \DecValTok{2}\NormalTok{, }\AttributeTok{fill =} \StringTok{"steelblue"}\NormalTok{, }\AttributeTok{color =} \StringTok{"white"}\NormalTok{)}
\end{Highlighting}
\end{Shaded}

\vspace{0.25em}

\begin{enumerate}
\def\labelenumi{(\alph{enumi})}
\item Which variable will be on the \textbf{x-axis} of the plot?

  \rule{2in}{0.4pt}
\item How wide is each bar (bin) in this histogram?

  \rule{2in}{0.4pt}
\item What color will the \textbf{inside} of the bars be?

  \rule{1.5in}{0.4pt}
\item What color will the \textbf{outline} of the bars be?

  \rule{1.5in}{0.4pt}
\end{enumerate}

\vspace{0.5em}

\subsection{9. (5 pts) Short answer --- debug the code}

A classmate runs this code and gets an error:

\begin{Shaded}
\begin{Highlighting}[]
\FunctionTok{ggplot}\NormalTok{(penguins, }\FunctionTok{aes}\NormalTok{(}\AttributeTok{x =}\NormalTok{ body\_mass\_g)) }\SpecialCharTok{\%\textgreater{}\%}
  \FunctionTok{geom\_histogram}\NormalTok{()}
\end{Highlighting}
\end{Shaded}

\vspace{0.25em}

\begin{enumerate}
\def\labelenumi{(\alph{enumi})}
\tightlist
\item Identify \textbf{what's wrong} with this code. (Hint: think about how \texttt{ggplot2} layers are joined.)
\end{enumerate}

\vspace{0.5em}

\rule{6.5in}{0.4pt}

\vspace{0.75em}

\rule{6.5in}{0.4pt}

\vspace{0.5em}

\begin{enumerate}
\def\labelenumi{(\alph{enumi})}
\setcounter{enumi}{1}
\tightlist
\item Write the \textbf{corrected} version of the code.
\end{enumerate}

\vspace{1em}

\rule{6.5in}{0.4pt}

\vspace{1em}

\rule{6.5in}{0.4pt}

\vspace{1em}

\vspace{0.5em}

\begin{center}
\textbf{End of practice quiz --- Total: 30 points}

\textit{Check your answers using the answer key.}
\end{center}

\end{document}